Section~6 of \cite{SharedDecoders2026} studied two-stage compositions $K_2 \circ K_1$ of linear
class-uniform shared decoders on $\Mat_{d\times d}(F^b)$ with $M_1$ and $M_2$ column classes, using
van der Waerden and Hales--Jewett arithmetic progressions to bound $d$ when $K_2 \circ K_1 = \tau$
and to show that $\rank_F(\tau - K_2 \circ K_1)/(bd) \to \infty$ as $d \to \infty$ at fixed
$M_1, M_2$ when $(d : F) \ne 0$. No multitape Turing machine is known to induce such a two-stage
linear class-uniform
factorization; working within the same algebraic model, this paper replaces arithmetic progressions
by Ramsey-free additive structures and proves two main composition bounds in Lean~4.

First (\Cref{thm:2a}), over any field $F$ and for $b, M_1, M_2 \ge 1$, replacing progressions by
affine Hilbert cubes extracted via iterated Cauchy--Schwarz proves that $K_2 \circ K_1 = \tau$ forces
\[
  d \;\le\; \max\bigl(8 M_1 M_2^2,\; (2 M_1)^{2^{M_2}}\bigr).
\]
This rules out exact two-stage linear factorizations for growing class schedules up to
$M_1 = o\bigl(d^{1/2^{M_2}}\bigr)$ at fixed $M_2$ (including $3$-AP-free Behrend partitions with
$M_1 = 2^{O(\sqrt{\log d})}$) and $M_2 \le (1-\eps)\log_2\log_2 d$ at fixed $M_1$ and
$\eps \in (0, 1)$ (giving $d \le \max(8M_1, 4M_1^2)$ at $M_2 = 1$), though at small
constant class counts such as $M_1 = M_2 = 2$ it gives $d \le 256$ where Corollary~6.4 of
\cite{SharedDecoders2026} gives $d \le 16$. Second (\Cref{thm:2b}), over any field
$F$ with $(d : F) \ne 0$, combining shifted $(2M_2)$-fold monochromatic sumsets, arbitrary-subset
Vandermonde pullbacks, sparse annihilating polynomials, and iterated K\H{o}v\'ari--S\'os--Tur\'an
extraction proves that for all $b, M_1, M_2 \ge 1$ and all $d \ge D_0(M_1, M_2)$,
\[
  \rank_F(\tau - K_2 \circ K_1)
  \;\ge\;
  c(M_1, M_2)\,b\,d\,(\log d)^{\frac{1}{2M_2 - 1}},
\]
giving $\Omega(b\,d\log d)$ when $M_2 = 1$ (\Cref{cor:2b-m2-one}), with explicit constants where
$D_0(1,1) = 16^{36} + 2 \approx 2.23 \times 10^{43}$. Finite-dimensional controls at $d = 13$,
$d = 4$, and $d = 12$ test the hypotheses of both bounds.

Three further sections prove one-stage linear shared-decoder and two-sided swap-code bounds on open
problems of \cite{SeventeenthPaper2026}, alongside two shorter companion sections:
(i)~on \cite[Problem~9.1(a)]{SeventeenthPaper2026}, exploiting decoder linearity on maximal
monochromatic cyclic runs proves an explicit single-shift intersection bound (\Cref{thm:1a-single-shift},
attained over every field on the alternating two-coloring for even $d = 2m \ge 2$ at $b = 1, s = 1$,
\Cref{thm:1a-sharpness}) and yields
$D_b(p) \le \frac{19}{32}\,bd^2 + \frac{3}{2}\,bd^{3/2}$ for two classes whenever
$\operatorname{char}(F) \nmid d$ (\Cref{thm:1b-two-class}, strictly below both the $\frac{41}{64}$
bound of \cite[Theorem~3.4]{SharedDecoders2026} and the $\frac{5}{8}$ nonlinear-decoder certificate
floor of \cite[Theorem~C.6]{SharedDecoders2026}) together with
$D_b(p) \le \bigl(1 - \frac{5}{4M} + \frac{1}{M^2} - \frac{1}{4M^3}\bigr)bd^2 + bMd$ for $M \ge 3$
at primes $d \ne \operatorname{char}(F)$ (\Cref{thm:1b-multi-class}, whose leading coefficient beats
the flag bound $1 - \frac{1}{2M}$ for every $M \ge 2$, while Theorem~3.5 of \cite{SharedDecoders2026}
has a smaller leading coefficient at $M = 50$ and $M \ge 55$, carries no $O(bMd)$ term, and applies to
composite $d$);
(ii)~on \cite[Problem~9.1(b)]{SeventeenthPaper2026}, coupled $\Z_6$ Fourier-mode elimination proves
$\dim_F(\mathcal{X} \cap \rot_1^{-1}(\mathcal{X})) \le 14$ and hence
$D_1\bigl(\{0,1,3\},\{2,4,5\}\bigr) = 25$ over every field with
$\operatorname{char}(F) \notin \{2, 3\}$ (\Cref{thm:aperiodic-six-25}, leaving characteristics $2$
and $3$ open);
(iii)~three-term swap codes $\sum_{t=0}^2 A_t X B_t^{\T} = X^{\T}$ over every field exceeding
$d(d+1)/2$ at $d \in \{6, 7, 9, 12\}$ and uniformly at $d = 4$ (\cite[Remark~5.6]{SeventeenthPaper2026}),
together with a characteristic-free two-sided coarse--fine bound on
\cite[Problem~9.3]{SeventeenthPaper2026} (\Cref{thm:two-sided-coarse-fine,cor:two-sided-win}) giving
leading coefficient $\frac{1}{2} + \frac{1}{4}\min\bigl(\beta(1+\alpha), \alpha(1+\beta)\bigr) <
\frac{1}{2} + \frac{1}{2}\min(\alpha, \beta)$ for all $\alpha, \beta \in (0,1)$ (and
$\frac{11}{16}d^2 + d < \frac{3}{4}d^2$ at $\alpha = \beta = 1/2$ for all even $d \ge 18$);
(iv)~an Operator Schmidt Rank pass lower bound ($t \ge \log_2 d$) for a restricted shuffle/unshuffle
pipeline paired with verified $O(L)$-time Turing machines in \texttt{lean-machine} for shuffle,
unshuffle, and within-row shear showing why $\OSR$ does not answer
\cite[Problem~9.5]{SeventeenthPaper2026}; and
(v)~cut-causal cone evaluation on frozen \leanname{Transpose.TM} fibres, a Lean formalization of the
$1/(\Delta+1)$ random-order bound of \cite[Proposition~4.14(a)--(b)]{FourthPaper2026}, and why
nonlinear set fibres and per-entry cut labels stop subspace and column-class arguments short of
\cite[Problem~9.4(b)]{SeventeenthPaper2026}. Statement~(T) remains open.
